\section{Related Work}\label{sec:related}

\subsection{Infrared and Visible Image Fusion}
Early deep fusion methods used auto-encoders and convolutional networks to learn source representations. DenseFuse~\cite{li2019densefuse} combines dense feature extraction with feature-level fusion, while SDNet~\cite{zhang2021sdnet} and U2Fusion~\cite{xu2022u2fusion} investigate learned representations and unsupervised objectives for multiple fusion tasks. FusionGAN~\cite{ma2019fusiongan} and TarDAL~\cite{liu2022tardal} introduce adversarial or target-aware supervision. PIAFusion~\cite{tang2022piafusion} adapts fusion to illumination, and SeAFusion~\cite{tang2022seafusion} connects image fusion with semantic perception. Their objectives reflect different priorities, including source reconstruction, target visibility, and semantic consistency.

Feature decomposition makes the organization of complementary information explicit. DIDFuse~\cite{zhao2020didfuse} learns background and detail features, and CDDFuse~\cite{zhao2023cddfuse} couples a Transformer-based base branch with an invertible neural network (INN) detail branch using cross-modal correlation constraints. More recently, FreeFusion~\cite{zhao2025freefusion} explores cross-reconstruction learning to reduce dependence on handcrafted fusion losses. Our work retains CDDFuse's two-stage decomposition framework and studies a compact architecture with additional feature decorrelation and image-domain regularization.

\subsection{Attention and Efficient Local Processing}
Attention provides content-dependent relationships for restoration and fusion. SwinFusion~\cite{ma2022swinfusion} combines intra-domain and inter-domain interactions through shifted-window attention. Restormer~\cite{zamir2022restormer} uses multi-DConv head transposed attention (MDTA) to estimate channel relationships from spatially aggregated features, avoiding the large pairwise spatial attention maps of conventional self-attention. FC-Former~\cite{wu2025fcformer} further investigates connections across attention modes for multi-source fusion. These studies motivate treating spatial interactions and channel relationships as complementary modeling choices.

Convolutional designs such as NAFNet~\cite{chen2022simple} show the value of efficient local operators in image restoration. Following this direction, our SGC Block combines multi-scale depthwise convolutions with gating and channel reweighting. We use it for shallow extraction and reconstruction, while reserving directional state-space processing and channel attention for the base branch. The focus is the allocation of operations across the fusion pipeline.

\subsection{State-Space Models for Visual Fusion}
Mamba~\cite{gu2023mamba} enables input-dependent sequence modeling with linear scaling in sequence length for fixed model dimensions. VMamba~\cite{liu2024vmamba} adapts scanning to visual data, and EfficientVMamba~\cite{pei2025efficientvmamba} explores efficient scan patterns. MambaIR~\cite{guo2024mambair} combines state-space processing with visual feature refinement, while MambaVision~\cite{hatamizadeh2025mambavision} studies a hybrid Mamba--Transformer backbone. Together, these works suggest combining spatial state propagation with local and channel-based processing.

For fusion, FusionMamba~\cite{xie2024fusionmamba} uses dynamic feature enhancement and cross-modal Mamba interaction. MambaDFuse~\cite{li2024mambadfuse} combines feature extraction with shallow and deep fusion. SCMFuse instead adapts hybrid spatial--channel modeling to a compact dual-branch decomposition network. Its scan path operates at reduced resolution, while an anchor and local path retain full-resolution features. DGF and MGF then use differences and magnitudes to guide feature aggregation. Our focus is how these operations can be integrated within a compact decomposition framework and evaluated through fusion quality and computational cost.
